\section{From the count estimate to full variation}\label{sec:completion}

We now assemble the smooth estimate and the endpoint estimate, remove
the restrictions on choices and inputs, and sum the count bounds to
obtain the full variation. No selection of a common partition for
different output points is made in this argument.

\begin{proof}[Proof of Theorem~\ref{thm:count}]
First let $F,G\in C_c^\infty(\R^2)$, use the pair maps
\eqref{eq:pair-maps}, and take $G_0\in C_c^\infty(\R^2)$.
For step choices of disjoint annuli and coefficients, the identity
\eqref{eq:hard-smooth} writes the hard linearization, multiplied by
$c_0\ne0$, as its smooth linearization minus its endpoint errors.
Propositions~\ref{prop:smooth-switches} and~\ref{prop:endpoint-errors}
therefore give
\begin{equation}\label{eq:hard-linearized-bound}
 \left|\iiint G_0(a,b)G_1(b,c)G_2(c,a)
     \sum_j\frac{\beta_j\mathbf1_{\{\varepsilon_j<|a+b+c|<R_j\}}}
                       {a+b+c}\dd a\dd b\dd c\right|
 \le C\sqrt n\log(2+n)\prod_v\|G_v\|_3.
\end{equation}
The interval list and coefficients in the sum may vary with $(a,b)$.

Fix a finite menu of such choices. Its measurable choice sets can be
approximated in measure on a compact set by finite unions of rectangles.
To see this, approximate each of the finitely many sets using Lebesgue
regularity and then refine all rectangle boundaries to one coordinate
partition. Resolve overlaps by a fixed ordering of the menu; the total
mismeasured set remains bounded by the sum of the approximation errors.
Assign any menu item on the remaining cells. For the fixed menu, the
annuli are uniformly bounded away from zero, and the integrals against
the smooth inputs are bounded on the compact support needed for $G_0$.
Consequently the corresponding linearized integrals converge.
Thus \eqref{eq:hard-linearized-bound} holds for arbitrary measurable
choices from a fixed finite menu.

Now fix a finite collection of lists of disjoint rational annuli, each
of length at most $n$. The maximum of their sums of absolute increments
is measurable. Choose a maximizing list measurably by resolving ties in
a fixed order. For each of its increments choose a coefficient from
$\{1,-1,i,-i\}$ whose product with that increment has real part at least
half its modulus. These choices are again measurable and range over
a finite menu. Testing \eqref{eq:hard-linearized-bound} against
nonnegative smooth compactly supported $G_0$ shows, by $L^3$ duality,
that this finite maximum has $L^{3/2}$ norm at most
$C\sqrt n\log(2+n)\|F\|_3\|G\|_3$.
The reflection $(x,y)=(-a,-b)$ preserves that norm.
Exhausting the countable family of all rational lists and applying
monotone convergence proves the count bound for smooth inputs.

For general complex $F,G\in L^3$, choose smooth compactly supported
approximants $F_m,G_m$. For every fixed annulus, bilinearity and
\eqref{eq:finite-annulus-bound} give convergence of
$B_{\varepsilon,R}(F_m,G_m)$ in $L^{3/2}$.
For a fixed finite collection of lists, the difference of the two
finite maxima is bounded by a finite sum of such truncation differences.
It therefore converges in $L^{3/2}$. Pass the estimate to the limit
for this finite maximum, then exhaust the rational lists once more.
Lemma~\ref{lem:finite-truncations} supplies the common representatives
required in the definition of $\mathcal S_n$.
\end{proof}

\begin{proof}[Proof of Theorem~\ref{thm:variation}]
For any finite rational partition, arrange the absolute annular
increments in nonincreasing order, denoting them by
$d_1\ge d_2\ge\cdots\ge0$. Whenever the index $2^j$ exists,
the corresponding largest increments are a permissible list in
$\mathcal S_{2^j}$, so
\[
 d_{2^j}\le2^{-j}\mathcal S_{2^j}(F,G).
\]
The group of indices $2^j\le l<2^{j+1}$ has at most $2^j$ members.
Bounding the full $\ell^r$ norm by the sum of the group norms yields,
uniformly in the chosen partition,
\begin{equation}\label{eq:rank-variation}
 V_r(F,G)\le\sum_{j\ge0}2^{j/r-j}\mathcal S_{2^j}(F,G)
 \quad\text{almost everywhere}.
\end{equation}
Theorem~\ref{thm:count} and Minkowski's inequality now imply
\[
 \|V_r(F,G)\|_{3/2}
 \le C\|F\|_3\|G\|_3
       \sum_{j\ge0}(1+j)2^{j(1/r-1/2)}.
\]
The series converges for every $r>2$. In particular, neither a bound
on the number of increments nor a restriction to dyadic endpoints
survives in the conclusion.
\end{proof}
